%% hopfion-framework/chemistry/icosahedrite.tex
%% Sources: main_paper11.tex
%% Icosahedrite Al63Cu24Fe13: phi^2 ratio, Fibonacci Fe=F7, quasicrystal
%% Auto-generated from project files. Edit source papers to update.


%════════════════════════════════════════════════════════════════════
% Icosahedrite Stoichiometry — Paper XI
%════════════════════════════════════════════════════════════════════
\begin{proposition}[Icosahedrite components on the spiral]
\label{P11:prop:icosahedrite}
The three components of icosahedrite have spiral indices:
\begin{equation}
  n_{\mathrm{Al}} = -0.8532, \quad
  n_{\mathrm{Fe}} = -0.5646, \quad
  n_{\mathrm{Cu}} = -0.5880,
  \label{P11:eq:icosahedrite_indices}
\end{equation}
giving the pairwise separations
\begin{equation}
  \Delta n(\mathrm{Fe,Cu}) = 0.0234, \quad
  \Delta n(\mathrm{Al,Cu}) = 0.265, \quad
  \Delta n(\mathrm{Al,Fe}) = 0.289.
  \label{P11:eq:icosahedrite_deltas}
\end{equation}
The stoichiometry $63:24:13$ satisfies
\begin{equation}
  \boxed{\frac{n_{\mathrm{Al}}}{n_{\mathrm{Cu}}} \;=\; \frac{63}{24}
  \;=\; 2.625 \;\approx\; \vp^2 \;=\; 2.618\,,\quad
  0.27\%\ \text{error},}
  \label{P11:eq:icosahedrite_phi2}
\end{equation}
and the residual count $\mathrm{Fe} = 100 - 63 - 24 = 13 = F_7$
is the seventh Fibonacci number.
\end{proposition}
\status{published}
\source{P11:prop:icosahedrite}

\begin{remark}[Three structural observations]
\label{P11:rem:icosahedrite_structure}
\textbf{(i) Fe--Cu near-degeneracy.}
The separation $\Delta n(\mathrm{Fe,Cu}) = 0.0234$ places Fe and Cu
in the same moderate-proximity cluster as Goldschmidt substitution pairs
(cf.\ the $\Delta n < 0.035$ threshold at which pairs achieve $>70\%$
recall).
Nickel, the primary geochemical substitute for both Fe and Cu,
has $n_{\mathrm{Ni}} = -0.5997$, giving
$\Delta n(\mathrm{Cu,Ni}) = 0.012$ and $\Delta n(\mathrm{Fe,Ni}) = 0.035$
--- the tightest cluster among the 3d metals.
This explains why natural icosahedrite samples show partial Ni/Fe
and Ni/Cu substitution~\cite{Bindi2009}, and why
Al--Cu--Fe--Ni is a common quasicrystal-forming system.

\textbf{(ii) Al plays a structurally distinct role.}
The large separation $\Delta n(\mathrm{Al,Fe}) \approx 0.29$
confirms that Al and the Fe/Cu cluster are in fundamentally different
spiral neighbourhoods.
Al ($n \approx -0.85$) sits between the group~I metals and the 3d cluster,
providing the structural (fcc-type packing) role in the quasicrystal
while Fe/Cu provide the electronic and magnetic character.
This spiral-based distinction is consistent with the known crystallography
of quasicrystals.

\textbf{(iii) Stoichiometry encodes $\vp^2 = \vp+1$.}
The relation $\mathrm{Al}/\mathrm{Cu} \approx \vp^2$ means the icosahedrite
stoichiometry reflects the fundamental golden-ratio identity $\vp^2 = \vp+1$
at $0.27\%$ accuracy.
Concretely: given $\mathrm{Cu} = 24 = 3F_6$, the spiral predicts
$\mathrm{Al} = \lfloor 24\vp^2 \rceil = \lfloor 62.83 \rceil = 63$,
and the complement $\mathrm{Fe} = 13 = F_7$ is a Fibonacci number.
The icosahedral symmetry group $2I$ is rooted in pentagon geometry
($k+2=5$ of the condensate); the quasicrystal whose symmetry group
is $2I$ has stoichiometry encoding $\vp^2$ to sub-percent accuracy.
\end{remark}
\status{published}
\source{P11:rem:icosahedrite_structure}

\begin{equation}
  \sigma_\mathrm{tot}
  = \sigma_\mathrm{Slater} + \sigma_\mathrm{xc}
  = 4.45 + 1.55
  = 6.00
  = Z - n_\mathrm{qn}
  = 8 - 2.
  \label{P11:eq:sigma_total}
\end{equation}

